% !TeX root = ../main.tex
\section{Related work}
\textbf{From predefined concepts to discovery.} TCAV measures prediction sensitivity along directions learned from examples of user-defined concepts~\cite{kim2018tcav}. Concept Bottleneck Models instead predict annotated concepts as an intermediate representation and permit interventions before the final prediction~\cite{koh2020cbm}. Both offer semantic interfaces, but rely on a supplied concept specification. Sparse dictionary learning removes this requirement by decomposing polysemantic activations into more monosemantic features without predefined concepts~\cite{bricken2023monosemanticity}, and Olah et al.~\cite{olah2020zoom} frame the resulting units and their connections as interpretable circuits rather than opaque activations. Building on this idea, Discover-then-Name applies SAE discovery before naming concepts and training downstream concept bottlenecks~\cite{rao2024discover}, while SpLiCE instead decomposes CLIP embeddings directly into sparse combinations of existing textual concepts, without training a new representation~\cite{bhalla2024splice}.

\textbf{VLMs and medical grounding.} Pach et al.~\cite{pach2025sae} evaluate SAE features in VLMs using a human-study benchmark and interventions, going beyond reconstruction alone. Closer to our clinical setting, Gong et al.~\cite{gong2025concepts} dissect neuron activations in medical image diagnostic models to build interpretable, concept-level explanations directly from opaque networks, motivating a neuron-to-concept alignment strategy for our chest radiograph setting rather than a supervised diagnostic head. MedConcept~\cite{haque2026medconcept} combines sparse decomposition of Merlin CT embeddings, an ontology-derived vocabulary, cosine-based naming, and frozen MedGemma~\cite{sellergren2026medgemma} judgments of concept--report consistency. We retain this separation between discovery and evaluation, while changing the image domain, backbone, vocabulary, and activation selection.

\textbf{Evaluating explanations.} Different evaluation methods capture distinct aspects of explanation quality: reconstruction measures how well embeddings are preserved, human inspection assesses semantic coherence, interventions test effects on model behavior, and report entailment assesses textual support. These approaches provide complementary evidence. Although LLM-based evaluation scales more readily than manual review, comparative judging studies have demonstrated susceptibility to position bias~\cite{shi2025judging}. We therefore use label permutation and independent validation to improve evaluation reliability, while recognizing that neither guarantees error-free judgments.